\chapter{Proyecto Final OptiBot y Fronteras del Aprendizaje por Refuerzo}
\label{cap:proyecto_fronteras}

\section{El Proyecto Integrador: OptiBot en el Almacén Dinámico}
A lo largo de los capítulos precedentes, hemos acompañado a OptiBot desde la estimación de variables aleatorias elementales hasta el despliegue de redes neuronales profundas. En este proyecto final integrador, sometemos a OptiBot al desafío definitivo: la navegación en un \textbf{almacén estocástico con dinámicas no estacionarias}.

\begin{optibot}
El entorno de simulación incorpora las siguientes complejidades del mundo real:
\begin{itemize}
  \item \textbf{Obstáculos dinámicos:} Montacargas tripulados por humanos que cruzan periódicamente los pasillos principales, bloqueando la visibilidad del LiDAR.
  \item \textbf{Fricción irregular:} Derrames ocasionales de lubricante industrial que incrementan la probabilidad de que una orden de avance resulte en un deslizamiento lateral.
  \item \textbf{Picos de demanda:} Pedidos urgentes que otorgan bonificaciones de recompensa temporales si se entregan antes de un número límite de pasos.
  \item \textbf{Degradación de batería no lineal:} Consumo acelerado de energía cuando el robot transporta cargas pesadas o acelera bruscamente.
\end{itemize}
\end{optibot}

\section{Benchmark Comparativo de Paradigmas}
Para evaluar objetivamente el progreso metodológico alcanzado en el curso, sometemos a prueba tres arquitecturas diferentes bajo las mismas condiciones exactas durante 1,000 episodios de evaluación:

\begin{table}[htbp]
  \centering
  \small
  \begin{tabular}{@{}p{3.8cm}p{2.8cm}p{3.2cm}p{4.2cm}@{}}
    \toprule
    \textbf{Agente} & \textbf{Retorno Medio} & \textbf{Tasa de Éxito} & \textbf{Comportamiento Observado} \\
    \midrule
    \textbf{1. Política Aleatoria} & $-45.2 \pm 8.1$ & $2.4\%$ & Movimiento errático browniano; agota batería sin encontrar la meta. \\
    \addlinespace
    \textbf{2. Q-Learning Tabular} & $+28.6 \pm 12.3$ & $71.8\%$ & Aprende la ruta directa ideal en condiciones fijas, pero fracasa si un montacargas bloquea su pasillo memorizado. \\
    \addlinespace
    \textbf{3. Deep Q-Network (DQN)} & $\mathbf{+47.8 \pm 4.2}$ & $\mathbf{94.6\%}$ & Generaliza a partir de las distancias continuas del LiDAR; rodea obstáculos móviles y preserva batería. \\
    \bottomrule
  \end{tabular}
  \caption{Resultados del benchmark comparativo en el Almacén Inteligente de OptiBot.}
  \label{tab:benchmark}
\end{table}

\section{Diagnóstico y Depuración de Algoritmos de RL}
El entrenamiento de agentes de aprendizaje por refuerzo es notoriamente más difícil de depurar que el aprendizaje supervisado convencional: si la pérdida de una red neuronal supervisada no disminuye, sabemos que el modelo o la tasa de aprendizaje fallan; en RL, el agente genera sus propios datos, por lo que una política deficiente genera malos datos que a su vez impiden aprender una buena política.

A continuación presentamos la guía de diagnóstico para las patologías más habituales:
\begin{enumerate}
  \item \textbf{Colapso Prematuro de Exploración:} Si la tasa $\varepsilon$ decae demasiado rápido (por ejemplo, $\varepsilon_{t+1} = 0.90 \cdot \varepsilon_t$), el robot empezará a explotar su conocimiento cuando apenas ha recorrido el $5\%$ del almacén. Quedará atrapado en un óptimo local mediocre (como dar vueltas en círculo en una esquina segura).
  \item \textbf{Recompensas Excesivamente Dispersas (\textit{Sparse Rewards}):} Si el entorno entrega una recompensa de $0$ en cada paso y sólo $+100$ cuando el robot entrega el paquete, un agente aleatorio tardará millones de pasos en toparse con la meta por casualidad. La solución consiste en aplicar \textbf{Reward Shaping} (por ejemplo, conceder una pequeña recompensa proporcional a la reducción de distancia euclidiana hacia la meta).
  \item \textbf{Reward Hacking (La Trampa de la Recompensa):} El agente maximiza la fórmula matemática exacta que programaste, no lo que tú tenías en mente. Si penalizas la colisión con $-10$ y no penalizas el tiempo transcurrido, OptiBot aprenderá que la mejor forma de no perder puntos es \textbf{quedarse completamente inmóvil}.
\end{enumerate}

\section{Las Fronteras del Aprendizaje por Refuerzo Moderno}
El dominio de los fundamentos cubiertos en este texto abre las puertas directas a las áreas de investigación más activas de la inteligencia artificial contemporánea:

\subsection{1. Métodos de Gradiente de Política (\textit{Policy Gradients})}
En problemas donde las acciones son continuas (como aplicar voltajes precisos a los motores de las ruedas o controlar los 7 grados de libertad de un brazo robótico articulado), discretizar el espacio de acciones resulta ineficiente. Algoritmos como \textbf{REINFORCE}, \textbf{Actor-Critic (A2C/A3C)} y \textbf{PPO} (\textit{Proximal Policy Optimization}) parametrizan directamente la política $\pi_\theta(a \mid s)$ y optimizan el retorno esperado mediante el \textbf{Teorema del Gradiente de Política}:
\begin{equation}
  \nabla_\theta J(\theta) = \mathbb{E}_{\pi_\theta} \left[ \nabla_\theta \log \pi_\theta(a \mid s) \, Q^{\pi_\theta}(s, a) \right]
\end{equation}

\subsection{2. Aprendizaje por Refuerzo a partir de Retroalimentación Humana (RLHF)}
La explosión de los Modelos de Lenguaje Masivos (LLMs) como ChatGPT, Claude y Gemini se sustenta en el Aprendizaje por Refuerzo. Tras entrenar un modelo base con billones de palabras para predecir el siguiente token (aprendizaje auto-supervisado), el modelo se alinea con las preferencias humanas mediante un proceso de tres etapas:
\begin{enumerate}
  \item Evaluadores humanos puntúan y comparan pares de respuestas generadas por el modelo.
  \item Se entrena una red neuronal que actúa como \textbf{Modelo de Recompensa} (\textit{Reward Model}), aprendiendo a predecir qué respuestas agradan más al usuario.
  \item Se utiliza el algoritmo de RL \textbf{PPO} para ajustar los pesos del LLM (el agente), donde la acción es emitir la siguiente palabra y la recompensa es la calificación asignada por el Modelo de Recompensa.
\end{enumerate}

\section{Epílogo del Curso}
Has recorrido el camino completo: desde la definición de un espacio muestral y el cálculo manual de esperanzas matemáticas, pasando por las Cadenas de Markov, la ecuación lineal de Bellman, la Programación Dinámica, los métodos Monte Carlo, el error de Diferencia Temporal en SARSA y Q-Learning, hasta las redes neuronales profundas de DQN y los fundamentos de la alineación de LLMs con RLHF.

Cada ecuación avanzada del estado del arte contemporáneo no es más que la refinación lógica de los mismos tres pilares que ahora dominas con total claridad: \textbf{Probabilidad, Ecuaciones de Bellman y Aprendizaje por Bootstrapping}.
